
\documentclass{article}
\usepackage{colortbl}
\usepackage{makecell}
\usepackage{multirow}
\usepackage{supertabular}

\begin{document}

\newcounter{utterance}

\centering \large Interaction Transcript for game `referencegame', experiment `random\_grids', episode 5 with reflexion{-}Qwen3.5{-}9B.
\vspace{24pt}

{ \footnotesize  \setcounter{utterance}{1}
\setlength{\tabcolsep}{0pt}
\begin{supertabular}{c@{$\;$}|p{.15\linewidth}@{}p{.15\linewidth}p{.15\linewidth}p{.15\linewidth}p{.15\linewidth}p{.15\linewidth}}
    \# & $\;$A & \multicolumn{4}{c}{Game Master} & $\;\:$B\\
    \hline

    \theutterance \stepcounter{utterance}  
    & & \multicolumn{4}{p{0.6\linewidth}}{
        \cellcolor[rgb]{0.9,0.9,0.9}{
            \makecell[{{p{\linewidth}}}]{
                \texttt{\tiny{[P1$\langle$GM]}}
                \texttt{You have attempted this same task before and did not succeed. These are the plans you wrote after those attempts.} \\
\texttt{Trial \#1: New plan: My previous description was too verbose and failed to create a unique, concise identifier by focusing on the specific pattern of empty cells (□) rather than listing every filled cell (X). To succeed, I will construct a shorter expression that highlights the distinctive "holes" in the target grid which differ from the distractors. Specifically, I will note that the target grid has empty cells at (1,3), (2,4), (3,2), (3,5), (4,1), (4,3), (5,2), and (5,3), forming a specific shape that distinguishes it from the distractors. I will phrase this as: "Expression: The 5x5 grid with empty cells at rows 1{-}3, 4, and 5 at columns 3, 4, 2, 5, 1, 3, 2, and 3 respectively, creating a unique sparse pattern." Wait, that's still clunky. Better: "Expression: The 5x5 grid where the empty cells form a specific asymmetric pattern, notably having two adjacent empty cells in the bottom row at columns 2 and 3, and a single empty cell at row 1, column 3, which distinguishes it from the distractors." Actually, let's be precise and concise: "Expression: The 5x5 grid with empty cells at (1,3), (2,4), (3,2), (3,5), (4,1), (4,3), (5,2), and (5,3)." This is direct and verifiable.} \\
\\ 
\texttt{Trial \#2: New plan: My previous attempt failed because I listed coordinates for empty cells that were either incorrect (e.g., (2,1) is filled in the target, (1,4) is filled) or redundant, leading to a description that did not uniquely distinguish the target from the distractors. Specifically, I missed the critical difference in the top row: the target has two adjacent empty cells at (1,3) and (1,4), whereas Distractor 1 has only one at (1,3) and Distractor 2 has none. I will revise the strategy to first identify the most distinctive feature—the pair of adjacent empty cells in the top row of the target grid—and construct the expression around this unique pattern, ensuring all listed coordinates are accurate.} \\
\\ 
\texttt{Expression: The 5x5 grid distinguished by having two adjacent empty cells in the top row at columns 3 and 4, along with empty cells at (2,1), (3,2), (3,5), (4,1), (4,3), (5,2), and (5,3).} \\
\\ 
\\ 
\texttt{You are given three grids, where each of them is 5 by 5 in size.} \\
\texttt{Grids have empty cells marked with "▢" and filled cells marked with "X".} \\
\texttt{Your task is to generate a referring expression that best describes the target grid while distinguishing it from the two other distractor grids.} \\
\texttt{The first grid is the target grid, and the following two grids are the distractors.} \\
\\ 
\texttt{Target grid:} \\
\\ 
\texttt{X X □ □ X} \\
\texttt{□ X X □ X} \\
\texttt{X □ X X □} \\
\texttt{□ X □ X X} \\
\texttt{X □ □ X X} \\
\\ 
\texttt{Distractor grid 1:} \\
\\ 
\texttt{X X □ X □} \\
\texttt{□ X X □ X} \\
\texttt{X □ □ X □} \\
\texttt{□ X □ □ X} \\
\texttt{X □ X □ X} \\
\\ 
\texttt{Distractor grid 2:} \\
\\ 
\texttt{□ X □ X □} \\
\texttt{X □ X □ X} \\
\texttt{□ X □ X □} \\
\texttt{X □ X □ X} \\
\texttt{□ X □ X □} \\
\\ 
\texttt{Instruction: Describe the target grid.} \\
\texttt{Generate the referring expression starting with the tag "Expression: " for the given target grid. Omit any other text.} \\
            }
        }
    }
    & \\ \\

    \theutterance \stepcounter{utterance}  
    & \multicolumn{4}{p{0.6\linewidth}}{
        \cellcolor[rgb]{0.9,0.9,0.9}{
            \makecell[{{p{\linewidth}}}]{
                \texttt{\tiny{[P1$\rangle$GM]}}
                \texttt{Expression: The 5x5 grid distinguished by having two adjacent empty cells in the top row at columns 3 and 4, along with empty cells at (2,1), (3,2), (3,5), (4,1), (4,3), (5,2), and (5,3).} \\
            }
        }
    }
    & & \\ \\

    \theutterance \stepcounter{utterance}  
    & & & \multicolumn{2}{p{0.3\linewidth}}{
        \cellcolor[rgb]{0.9,0.9,0.9}{
            \makecell[{{p{\linewidth}}}]{
                \texttt{\tiny{[GM$|$GM]}}
                \texttt{Expression: The 5x5 grid distinguished by having two adjacent empty cells in the top row at columns 3 and 4, along with empty cells at (2,1), (3,2), (3,5), (4,1), (4,3), (5,2), and (5,3).} \\
            }
        }
    }
    & & \\ \\

    \theutterance \stepcounter{utterance}  
    & & \multicolumn{4}{p{0.6\linewidth}}{
        \cellcolor[rgb]{0.9,0.9,0.9}{
            \makecell[{{p{\linewidth}}}]{
                \texttt{\tiny{[GM$\rangle$P2]}}
                \texttt{You have attempted this same task before and did not succeed. These are the plans you wrote after those attempts.} \\
\texttt{Trial \#1: New plan: I incorrectly identified the "First" grid because I likely misaligned the row descriptions or miscounted the filled cells. Specifically, the expression requires the first row to have filled cells at positions 1 and 5, but the First grid has them at 1 and 3. The Third grid matches the description perfectly: Row 1 (1,5), Row 2 (2,3,5), Row 3 (1,3,4), Row 4 (2,4), and Row 5 (1,4,5). Next time, I will systematically verify each row of every grid against the expression's specific coordinate requirements before making a selection, ensuring the Third grid is chosen.} \\
\texttt{Answer: Third} \\
\\ 
\texttt{Trial \#2: New plan: I failed because I did not rigorously verify the specific coordinates of the empty cells ("▢") against the provided expression. The expression explicitly lists empty cells at (1,3), (1,4), (2,1), (2,4), (3,2), (3,5), (4,1), (4,3), (5,2), and (5,3). I will now systematically check each grid by mapping these exact coordinates:} \\
\texttt{1.  **Check Grid 1**: Row 1 has empty cells at (1,1), (1,3), (1,5). This does not match the required (1,3) and (1,4).} \\
\texttt{2.  **Check Grid 2**: Row 1 has empty cells at (1,3), (1,5). This does not match the required (1,3) and (1,4).} \\
\texttt{3.  **Check Grid 3**: Row 1 has empty cells at (1,3), (1,4). Row 2 has empty cells at (2,1), (2,4). Row 3 has empty cells at (3,2), (3,5). Row 4 has empty cells at (4,1), (4,3). Row 5 has empty cells at (5,2), (5,3). This matches the expression perfectly.} \\
\texttt{Next time, I will strictly cross{-}reference every coordinate in the expression list with the visual grid before answering.} \\
\texttt{Answer: Third} \\
\\ 
\\ 
\texttt{You are given three grids, where each of them is 5 by 5 in size.} \\
\texttt{Grids have empty cells marked with "▢" and filled cells marked with "X".} \\
\texttt{You are also given a referring expression that describes one of the given grids.} \\
\texttt{Your task is to select the grid that matches the given referring expression.} \\
\texttt{Generate only the number (in text) of the grid that matches the given expression by selecting first, second, or third.} \\
\\ 
\texttt{First:} \\
\\ 
\texttt{□ X □ X □} \\
\texttt{X □ X □ X} \\
\texttt{□ X □ X □} \\
\texttt{X □ X □ X} \\
\texttt{□ X □ X □} \\
\\ 
\texttt{Second:} \\
\\ 
\texttt{X X □ X □} \\
\texttt{□ X X □ X} \\
\texttt{X □ □ X □} \\
\texttt{□ X □ □ X} \\
\texttt{X □ X □ X} \\
\\ 
\texttt{Third:} \\
\\ 
\texttt{X X □ □ X} \\
\texttt{□ X X □ X} \\
\texttt{X □ X X □} \\
\texttt{□ X □ X X} \\
\texttt{X □ □ X X} \\
\\ 
\texttt{Expression: The 5x5 grid distinguished by having two adjacent empty cells in the top row at columns 3 and 4, along with empty cells at (2,1), (3,2), (3,5), (4,1), (4,3), (5,2), and (5,3).} \\
\texttt{Question: Which grid does the expression refer to?} \\
\texttt{Start with the tag "Answer: ", followed by your selection. Omit any other text.} \\
            }
        }
    }
    & \\ \\

    \theutterance \stepcounter{utterance}  
    & & & \multicolumn{4}{p{0.6\linewidth}}{
        \cellcolor[rgb]{0.9,0.9,0.9}{
            \makecell[{{p{\linewidth}}}]{
                \texttt{\tiny{[GM$\langle$P2]}}
                \texttt{Answer: Third} \\
            }
        }
    }
     \\ \\

    \theutterance \stepcounter{utterance}  
    & & & \multicolumn{2}{p{0.3\linewidth}}{
        \cellcolor[rgb]{0.9,0.9,0.9}{
            \makecell[{{p{\linewidth}}}]{
                \texttt{\tiny{[GM$|$GM]}}
                \texttt{third} \\
            }
        }
    }
    & & \\ \\

\end{supertabular}
}

\end{document}
